\subsection{Distributed Systems and Reliability}
\paragraph{MCQ-1457 [medium]} A design requires this behavior: Under a network partition, a distributed system must trade availability against consistent reads and writes. Which mechanism should the team study?
\begin{enumerate}[label=\Alph*.]
\item Queues and backpressure
\item CAP and consistency
\item Caching
\item Retries and timeouts
\end{enumerate}
\noindent\textbf{Answer.} CAP and consistency
\par\textit{Under a network partition, a distributed system must trade availability against consistent reads and writes. Tradeoff: Stronger guarantees simplify application invariants but increase latency and unavailability. Watch for: Quoting CAP without identifying partitions and operations does not guide design.}
\paragraph{MCQ-1458 [hard]} The team proposes CAP and consistency for a real workload. Which finding gives the correct decision boundary?
\begin{enumerate}[label=\Alph*.]
\item Success rate improves, while retries amplify overload and tail latency.
\item Stronger guarantees simplify application invariants but increase latency and unavailability.
\item Smoothing improves utilization, while waiting increases latency and stale work.
\item Hit rates improve cost, but correctness depends on keys, TTLs, and invalidation.
\end{enumerate}
\noindent\textbf{Answer.} Stronger guarantees simplify application invariants but increase latency and unavailability.
\par\textit{Under a network partition, a distributed system must trade availability against consistent reads and writes. Tradeoff: Stronger guarantees simplify application invariants but increase latency and unavailability. Watch for: Quoting CAP without identifying partitions and operations does not guide design.}
\paragraph{MCQ-1459 [hard]} Before releasing CAP and consistency, which failure deserves a dedicated regression test?
\begin{enumerate}[label=\Alph*.]
\item Layered automatic retries create multiplicative request storms.
\item Omitting tenant, model, prompt version, or authorization from a cache key leaks data.
\item An unbounded queue converts overload into memory exhaustion.
\item Quoting CAP without identifying partitions and operations does not guide design.
\end{enumerate}
\noindent\textbf{Answer.} Quoting CAP without identifying partitions and operations does not guide design.
\par\textit{Under a network partition, a distributed system must trade availability against consistent reads and writes. Tradeoff: Stronger guarantees simplify application invariants but increase latency and unavailability. Watch for: Quoting CAP without identifying partitions and operations does not guide design.}
\paragraph{MCQ-1460 [medium]} A design review reveals this mistake: Quoting CAP without identifying partitions and operations does not guide design. Which mechanism should be traced first?
\begin{enumerate}[label=\Alph*.]
\item Caching
\item Queues and backpressure
\item CAP and consistency
\item Retries and timeouts
\end{enumerate}
\noindent\textbf{Answer.} CAP and consistency
\par\textit{Under a network partition, a distributed system must trade availability against consistent reads and writes. Tradeoff: Stronger guarantees simplify application invariants but increase latency and unavailability. Watch for: Quoting CAP without identifying partitions and operations does not guide design.}
\paragraph{MCQ-1461 [hard]} Which causal explanation of CAP and consistency connects what it does to when it is worth its cost?
\begin{enumerate}[label=\Alph*.]
\item Mechanism: Timeouts bound uncertainty; retries use budgets, backoff, jitter, and idempotency for transient failures. Consequence: Success rate improves, while retries amplify overload and tail latency.
\item Mechanism: Caches trade staleness and invalidation complexity for lower latency and load. Consequence: Hit rates improve cost, but correctness depends on keys, TTLs, and invalidation.
\item Mechanism: Bounded queues absorb bursts while admission control and backpressure prevent producers from overwhelming consumers. Consequence: Smoothing improves utilization, while waiting increases latency and stale work.
\item Mechanism: Under a network partition, a distributed system must trade availability against consistent reads and writes. Consequence: Stronger guarantees simplify application invariants but increase latency and unavailability.
\end{enumerate}
\noindent\textbf{Answer.} Mechanism: Under a network partition, a distributed system must trade availability against consistent reads and writes. Consequence: Stronger guarantees simplify application invariants but increase latency and unavailability.
\par\textit{Under a network partition, a distributed system must trade availability against consistent reads and writes. Tradeoff: Stronger guarantees simplify application invariants but increase latency and unavailability. Watch for: Quoting CAP without identifying partitions and operations does not guide design.}
\paragraph{FC-1462 [medium]} Define CAP and consistency in interview-ready language.
\noindent\textbf{Answer.} Under a network partition, a distributed system must trade availability against consistent reads and writes.
\par\textit{Under a network partition, a distributed system must trade availability against consistent reads and writes.}
\paragraph{FC-1463 [hard]} State the main tradeoff and production pitfall for CAP and consistency.
\noindent\textbf{Answer.} Tradeoff: Stronger guarantees simplify application invariants but increase latency and unavailability. Pitfall: Quoting CAP without identifying partitions and operations does not guide design.
\par\textit{Tradeoff: Stronger guarantees simplify application invariants but increase latency and unavailability. Pitfall: Quoting CAP without identifying partitions and operations does not guide design.}
\paragraph{LA-1464 [hard]} Explain CAP and consistency from first principles, then show how you would evaluate and operate it in production.
\noindent\textbf{Answer.} Start with the mechanism: Under a network partition, a distributed system must trade availability against consistent reads and writes. Make the tradeoff explicit: Stronger guarantees simplify application invariants but increase latency and unavailability. Define workload-specific offline metrics and a latency/cost budget, test important slices, instrument the critical path, and create a rollback criterion. The production review must explicitly guard against this failure: Quoting CAP without identifying partitions and operations does not guide design.
\par\textit{A strong answer connects mechanism, measurable tradeoffs, failure isolation, observability, and rollback rather than listing product features.}
\paragraph{MCQ-1465 [medium]} A design requires this behavior: Protocols such as Raft replicate an ordered log so nodes agree despite failures. Which mechanism should the team study?
\begin{enumerate}[label=\Alph*.]
\item Exactly-once effects
\item Capacity planning
\item Consensus
\item CAP and consistency
\end{enumerate}
\noindent\textbf{Answer.} Consensus
\par\textit{Protocols such as Raft replicate an ordered log so nodes agree despite failures. Tradeoff: Coordination provides a consistent control plane but costs extra messages and quorum latency. Watch for: Deploying a quorum across unreliable links can reduce availability.}
\paragraph{MCQ-1466 [hard]} The team proposes Consensus for a real workload. Which finding gives the correct decision boundary?
\begin{enumerate}[label=\Alph*.]
\item Headroom handles bursts but idle accelerators are expensive.
\item Coordination provides a consistent control plane but costs extra messages and quorum latency.
\item Stronger guarantees simplify application invariants but increase latency and unavailability.
\item Business effects can be effectively once, but state and retention are needed.
\end{enumerate}
\noindent\textbf{Answer.} Coordination provides a consistent control plane but costs extra messages and quorum latency.
\par\textit{Protocols such as Raft replicate an ordered log so nodes agree despite failures. Tradeoff: Coordination provides a consistent control plane but costs extra messages and quorum latency. Watch for: Deploying a quorum across unreliable links can reduce availability.}
\paragraph{MCQ-1467 [hard]} Before releasing Consensus, which failure deserves a dedicated regression test?
\begin{enumerate}[label=\Alph*.]
\item Quoting CAP without identifying partitions and operations does not guide design.
\item Equating broker delivery semantics with end-to-end exactly-once effects is incorrect.
\item Deploying a quorum across unreliable links can reduce availability.
\item Planning with average sequence length underestimates tail memory and latency.
\end{enumerate}
\noindent\textbf{Answer.} Deploying a quorum across unreliable links can reduce availability.
\par\textit{Protocols such as Raft replicate an ordered log so nodes agree despite failures. Tradeoff: Coordination provides a consistent control plane but costs extra messages and quorum latency. Watch for: Deploying a quorum across unreliable links can reduce availability.}
\paragraph{MCQ-1468 [medium]} A design review reveals this mistake: Deploying a quorum across unreliable links can reduce availability. Which mechanism should be traced first?
\begin{enumerate}[label=\Alph*.]
\item Capacity planning
\item Exactly-once effects
\item CAP and consistency
\item Consensus
\end{enumerate}
\noindent\textbf{Answer.} Consensus
\par\textit{Protocols such as Raft replicate an ordered log so nodes agree despite failures. Tradeoff: Coordination provides a consistent control plane but costs extra messages and quorum latency. Watch for: Deploying a quorum across unreliable links can reduce availability.}
\paragraph{MCQ-1469 [hard]} Which causal explanation of Consensus connects what it does to when it is worth its cost?
\begin{enumerate}[label=\Alph*.]
\item Mechanism: Workload models connect arrival rate, prompt and output tokens, service time, concurrency, memory, and SLO targets. Consequence: Headroom handles bursts but idle accelerators are expensive.
\item Mechanism: Under a network partition, a distributed system must trade availability against consistent reads and writes. Consequence: Stronger guarantees simplify application invariants but increase latency and unavailability.
\item Mechanism: Protocols such as Raft replicate an ordered log so nodes agree despite failures. Consequence: Coordination provides a consistent control plane but costs extra messages and quorum latency.
\item Mechanism: Systems typically combine at-least-once delivery with idempotent consumers, dedupe keys, and transactional outboxes. Consequence: Business effects can be effectively once, but state and retention are needed.
\end{enumerate}
\noindent\textbf{Answer.} Mechanism: Protocols such as Raft replicate an ordered log so nodes agree despite failures. Consequence: Coordination provides a consistent control plane but costs extra messages and quorum latency.
\par\textit{Protocols such as Raft replicate an ordered log so nodes agree despite failures. Tradeoff: Coordination provides a consistent control plane but costs extra messages and quorum latency. Watch for: Deploying a quorum across unreliable links can reduce availability.}
\paragraph{FC-1470 [medium]} Define Consensus in interview-ready language.
\noindent\textbf{Answer.} Protocols such as Raft replicate an ordered log so nodes agree despite failures.
\par\textit{Protocols such as Raft replicate an ordered log so nodes agree despite failures.}
\paragraph{FC-1471 [hard]} State the main tradeoff and production pitfall for Consensus.
\noindent\textbf{Answer.} Tradeoff: Coordination provides a consistent control plane but costs extra messages and quorum latency. Pitfall: Deploying a quorum across unreliable links can reduce availability.
\par\textit{Tradeoff: Coordination provides a consistent control plane but costs extra messages and quorum latency. Pitfall: Deploying a quorum across unreliable links can reduce availability.}
\paragraph{LA-1472 [hard]} Explain Consensus from first principles, then show how you would evaluate and operate it in production.
\noindent\textbf{Answer.} Start with the mechanism: Protocols such as Raft replicate an ordered log so nodes agree despite failures. Make the tradeoff explicit: Coordination provides a consistent control plane but costs extra messages and quorum latency. Define workload-specific offline metrics and a latency/cost budget, test important slices, instrument the critical path, and create a rollback criterion. The production review must explicitly guard against this failure: Deploying a quorum across unreliable links can reduce availability.
\par\textit{A strong answer connects mechanism, measurable tradeoffs, failure isolation, observability, and rollback rather than listing product features.}
\paragraph{MCQ-1473 [medium]} A design requires this behavior: Bounded queues absorb bursts while admission control and backpressure prevent producers from overwhelming consumers. Which mechanism should the team study?
\begin{enumerate}[label=\Alph*.]
\item CAP and consistency
\item Consensus
\item Caching
\item Queues and backpressure
\end{enumerate}
\noindent\textbf{Answer.} Queues and backpressure
\par\textit{Bounded queues absorb bursts while admission control and backpressure prevent producers from overwhelming consumers. Tradeoff: Smoothing improves utilization, while waiting increases latency and stale work. Watch for: An unbounded queue converts overload into memory exhaustion.}
\paragraph{MCQ-1474 [hard]} The team proposes Queues and backpressure for a real workload. Which finding gives the correct decision boundary?
\begin{enumerate}[label=\Alph*.]
\item Stronger guarantees simplify application invariants but increase latency and unavailability.
\item Smoothing improves utilization, while waiting increases latency and stale work.
\item Coordination provides a consistent control plane but costs extra messages and quorum latency.
\item Hit rates improve cost, but correctness depends on keys, TTLs, and invalidation.
\end{enumerate}
\noindent\textbf{Answer.} Smoothing improves utilization, while waiting increases latency and stale work.
\par\textit{Bounded queues absorb bursts while admission control and backpressure prevent producers from overwhelming consumers. Tradeoff: Smoothing improves utilization, while waiting increases latency and stale work. Watch for: An unbounded queue converts overload into memory exhaustion.}
\paragraph{MCQ-1475 [hard]} Before releasing Queues and backpressure, which failure deserves a dedicated regression test?
\begin{enumerate}[label=\Alph*.]
\item An unbounded queue converts overload into memory exhaustion.
\item Deploying a quorum across unreliable links can reduce availability.
\item Quoting CAP without identifying partitions and operations does not guide design.
\item Omitting tenant, model, prompt version, or authorization from a cache key leaks data.
\end{enumerate}
\noindent\textbf{Answer.} An unbounded queue converts overload into memory exhaustion.
\par\textit{Bounded queues absorb bursts while admission control and backpressure prevent producers from overwhelming consumers. Tradeoff: Smoothing improves utilization, while waiting increases latency and stale work. Watch for: An unbounded queue converts overload into memory exhaustion.}
\paragraph{MCQ-1476 [medium]} A design review reveals this mistake: An unbounded queue converts overload into memory exhaustion. Which mechanism should be traced first?
\begin{enumerate}[label=\Alph*.]
\item CAP and consistency
\item Queues and backpressure
\item Consensus
\item Caching
\end{enumerate}
\noindent\textbf{Answer.} Queues and backpressure
\par\textit{Bounded queues absorb bursts while admission control and backpressure prevent producers from overwhelming consumers. Tradeoff: Smoothing improves utilization, while waiting increases latency and stale work. Watch for: An unbounded queue converts overload into memory exhaustion.}
\paragraph{MCQ-1477 [hard]} Which causal explanation of Queues and backpressure connects what it does to when it is worth its cost?
\begin{enumerate}[label=\Alph*.]
\item Mechanism: Bounded queues absorb bursts while admission control and backpressure prevent producers from overwhelming consumers. Consequence: Smoothing improves utilization, while waiting increases latency and stale work.
\item Mechanism: Under a network partition, a distributed system must trade availability against consistent reads and writes. Consequence: Stronger guarantees simplify application invariants but increase latency and unavailability.
\item Mechanism: Protocols such as Raft replicate an ordered log so nodes agree despite failures. Consequence: Coordination provides a consistent control plane but costs extra messages and quorum latency.
\item Mechanism: Caches trade staleness and invalidation complexity for lower latency and load. Consequence: Hit rates improve cost, but correctness depends on keys, TTLs, and invalidation.
\end{enumerate}
\noindent\textbf{Answer.} Mechanism: Bounded queues absorb bursts while admission control and backpressure prevent producers from overwhelming consumers. Consequence: Smoothing improves utilization, while waiting increases latency and stale work.
\par\textit{Bounded queues absorb bursts while admission control and backpressure prevent producers from overwhelming consumers. Tradeoff: Smoothing improves utilization, while waiting increases latency and stale work. Watch for: An unbounded queue converts overload into memory exhaustion.}
\paragraph{FC-1478 [medium]} Define Queues and backpressure in interview-ready language.
\noindent\textbf{Answer.} Bounded queues absorb bursts while admission control and backpressure prevent producers from overwhelming consumers.
\par\textit{Bounded queues absorb bursts while admission control and backpressure prevent producers from overwhelming consumers.}
\paragraph{FC-1479 [hard]} State the main tradeoff and production pitfall for Queues and backpressure.
\noindent\textbf{Answer.} Tradeoff: Smoothing improves utilization, while waiting increases latency and stale work. Pitfall: An unbounded queue converts overload into memory exhaustion.
\par\textit{Tradeoff: Smoothing improves utilization, while waiting increases latency and stale work. Pitfall: An unbounded queue converts overload into memory exhaustion.}
\paragraph{LA-1480 [hard]} Explain Queues and backpressure from first principles, then show how you would evaluate and operate it in production.
\noindent\textbf{Answer.} Start with the mechanism: Bounded queues absorb bursts while admission control and backpressure prevent producers from overwhelming consumers. Make the tradeoff explicit: Smoothing improves utilization, while waiting increases latency and stale work. Define workload-specific offline metrics and a latency/cost budget, test important slices, instrument the critical path, and create a rollback criterion. The production review must explicitly guard against this failure: An unbounded queue converts overload into memory exhaustion.
\par\textit{A strong answer connects mechanism, measurable tradeoffs, failure isolation, observability, and rollback rather than listing product features.}
\paragraph{MCQ-1481 [medium]} A design requires this behavior: Timeouts bound uncertainty; retries use budgets, backoff, jitter, and idempotency for transient failures. Which mechanism should the team study?
\begin{enumerate}[label=\Alph*.]
\item Retries and timeouts
\item CAP and consistency
\item Queues and backpressure
\item Multi-region design
\end{enumerate}
\noindent\textbf{Answer.} Retries and timeouts
\par\textit{Timeouts bound uncertainty; retries use budgets, backoff, jitter, and idempotency for transient failures. Tradeoff: Success rate improves, while retries amplify overload and tail latency. Watch for: Layered automatic retries create multiplicative request storms.}
\paragraph{MCQ-1482 [hard]} The team proposes Retries and timeouts for a real workload. Which finding gives the correct decision boundary?
\begin{enumerate}[label=\Alph*.]
\item Success rate improves, while retries amplify overload and tail latency.
\item Smoothing improves utilization, while waiting increases latency and stale work.
\item Active-active improves availability and latency but makes writes and conflict resolution harder.
\item Stronger guarantees simplify application invariants but increase latency and unavailability.
\end{enumerate}
\noindent\textbf{Answer.} Success rate improves, while retries amplify overload and tail latency.
\par\textit{Timeouts bound uncertainty; retries use budgets, backoff, jitter, and idempotency for transient failures. Tradeoff: Success rate improves, while retries amplify overload and tail latency. Watch for: Layered automatic retries create multiplicative request storms.}
\paragraph{MCQ-1483 [hard]} Before releasing Retries and timeouts, which failure deserves a dedicated regression test?
\begin{enumerate}[label=\Alph*.]
\item An unbounded queue converts overload into memory exhaustion.
\item A standby region that never receives realistic traffic may fail during disaster recovery.
\item Quoting CAP without identifying partitions and operations does not guide design.
\item Layered automatic retries create multiplicative request storms.
\end{enumerate}
\noindent\textbf{Answer.} Layered automatic retries create multiplicative request storms.
\par\textit{Timeouts bound uncertainty; retries use budgets, backoff, jitter, and idempotency for transient failures. Tradeoff: Success rate improves, while retries amplify overload and tail latency. Watch for: Layered automatic retries create multiplicative request storms.}
\paragraph{MCQ-1484 [medium]} A design review reveals this mistake: Layered automatic retries create multiplicative request storms. Which mechanism should be traced first?
\begin{enumerate}[label=\Alph*.]
\item Queues and backpressure
\item CAP and consistency
\item Retries and timeouts
\item Multi-region design
\end{enumerate}
\noindent\textbf{Answer.} Retries and timeouts
\par\textit{Timeouts bound uncertainty; retries use budgets, backoff, jitter, and idempotency for transient failures. Tradeoff: Success rate improves, while retries amplify overload and tail latency. Watch for: Layered automatic retries create multiplicative request storms.}
\paragraph{MCQ-1485 [hard]} Which causal explanation of Retries and timeouts connects what it does to when it is worth its cost?
\begin{enumerate}[label=\Alph*.]
\item Mechanism: Regions trade user latency, residency, failover speed, replication lag, and operational complexity. Consequence: Active-active improves availability and latency but makes writes and conflict resolution harder.
\item Mechanism: Timeouts bound uncertainty; retries use budgets, backoff, jitter, and idempotency for transient failures. Consequence: Success rate improves, while retries amplify overload and tail latency.
\item Mechanism: Bounded queues absorb bursts while admission control and backpressure prevent producers from overwhelming consumers. Consequence: Smoothing improves utilization, while waiting increases latency and stale work.
\item Mechanism: Under a network partition, a distributed system must trade availability against consistent reads and writes. Consequence: Stronger guarantees simplify application invariants but increase latency and unavailability.
\end{enumerate}
\noindent\textbf{Answer.} Mechanism: Timeouts bound uncertainty; retries use budgets, backoff, jitter, and idempotency for transient failures. Consequence: Success rate improves, while retries amplify overload and tail latency.
\par\textit{Timeouts bound uncertainty; retries use budgets, backoff, jitter, and idempotency for transient failures. Tradeoff: Success rate improves, while retries amplify overload and tail latency. Watch for: Layered automatic retries create multiplicative request storms.}
\paragraph{FC-1486 [medium]} Define Retries and timeouts in interview-ready language.
\noindent\textbf{Answer.} Timeouts bound uncertainty; retries use budgets, backoff, jitter, and idempotency for transient failures.
\par\textit{Timeouts bound uncertainty; retries use budgets, backoff, jitter, and idempotency for transient failures.}
\paragraph{FC-1487 [hard]} State the main tradeoff and production pitfall for Retries and timeouts.
\noindent\textbf{Answer.} Tradeoff: Success rate improves, while retries amplify overload and tail latency. Pitfall: Layered automatic retries create multiplicative request storms.
\par\textit{Tradeoff: Success rate improves, while retries amplify overload and tail latency. Pitfall: Layered automatic retries create multiplicative request storms.}
\paragraph{LA-1488 [hard]} Explain Retries and timeouts from first principles, then show how you would evaluate and operate it in production.
\noindent\textbf{Answer.} Start with the mechanism: Timeouts bound uncertainty; retries use budgets, backoff, jitter, and idempotency for transient failures. Make the tradeoff explicit: Success rate improves, while retries amplify overload and tail latency. Define workload-specific offline metrics and a latency/cost budget, test important slices, instrument the critical path, and create a rollback criterion. The production review must explicitly guard against this failure: Layered automatic retries create multiplicative request storms.
\par\textit{A strong answer connects mechanism, measurable tradeoffs, failure isolation, observability, and rollback rather than listing product features.}
\paragraph{MCQ-1489 [medium]} A design requires this behavior: Caches trade staleness and invalidation complexity for lower latency and load. Which mechanism should the team study?
\begin{enumerate}[label=\Alph*.]
\item Caching
\item Retries and timeouts
\item Capacity planning
\item CAP and consistency
\end{enumerate}
\noindent\textbf{Answer.} Caching
\par\textit{Caches trade staleness and invalidation complexity for lower latency and load. Tradeoff: Hit rates improve cost, but correctness depends on keys, TTLs, and invalidation. Watch for: Omitting tenant, model, prompt version, or authorization from a cache key leaks data.}
\paragraph{MCQ-1490 [hard]} The team proposes Caching for a real workload. Which finding gives the correct decision boundary?
\begin{enumerate}[label=\Alph*.]
\item Hit rates improve cost, but correctness depends on keys, TTLs, and invalidation.
\item Headroom handles bursts but idle accelerators are expensive.
\item Success rate improves, while retries amplify overload and tail latency.
\item Stronger guarantees simplify application invariants but increase latency and unavailability.
\end{enumerate}
\noindent\textbf{Answer.} Hit rates improve cost, but correctness depends on keys, TTLs, and invalidation.
\par\textit{Caches trade staleness and invalidation complexity for lower latency and load. Tradeoff: Hit rates improve cost, but correctness depends on keys, TTLs, and invalidation. Watch for: Omitting tenant, model, prompt version, or authorization from a cache key leaks data.}
\paragraph{MCQ-1491 [hard]} Before releasing Caching, which failure deserves a dedicated regression test?
\begin{enumerate}[label=\Alph*.]
\item Quoting CAP without identifying partitions and operations does not guide design.
\item Omitting tenant, model, prompt version, or authorization from a cache key leaks data.
\item Planning with average sequence length underestimates tail memory and latency.
\item Layered automatic retries create multiplicative request storms.
\end{enumerate}
\noindent\textbf{Answer.} Omitting tenant, model, prompt version, or authorization from a cache key leaks data.
\par\textit{Caches trade staleness and invalidation complexity for lower latency and load. Tradeoff: Hit rates improve cost, but correctness depends on keys, TTLs, and invalidation. Watch for: Omitting tenant, model, prompt version, or authorization from a cache key leaks data.}
\paragraph{MCQ-1492 [medium]} A design review reveals this mistake: Omitting tenant, model, prompt version, or authorization from a cache key leaks data. Which mechanism should be traced first?
\begin{enumerate}[label=\Alph*.]
\item CAP and consistency
\item Capacity planning
\item Retries and timeouts
\item Caching
\end{enumerate}
\noindent\textbf{Answer.} Caching
\par\textit{Caches trade staleness and invalidation complexity for lower latency and load. Tradeoff: Hit rates improve cost, but correctness depends on keys, TTLs, and invalidation. Watch for: Omitting tenant, model, prompt version, or authorization from a cache key leaks data.}
\paragraph{MCQ-1493 [hard]} Which causal explanation of Caching connects what it does to when it is worth its cost?
\begin{enumerate}[label=\Alph*.]
\item Mechanism: Under a network partition, a distributed system must trade availability against consistent reads and writes. Consequence: Stronger guarantees simplify application invariants but increase latency and unavailability.
\item Mechanism: Workload models connect arrival rate, prompt and output tokens, service time, concurrency, memory, and SLO targets. Consequence: Headroom handles bursts but idle accelerators are expensive.
\item Mechanism: Timeouts bound uncertainty; retries use budgets, backoff, jitter, and idempotency for transient failures. Consequence: Success rate improves, while retries amplify overload and tail latency.
\item Mechanism: Caches trade staleness and invalidation complexity for lower latency and load. Consequence: Hit rates improve cost, but correctness depends on keys, TTLs, and invalidation.
\end{enumerate}
\noindent\textbf{Answer.} Mechanism: Caches trade staleness and invalidation complexity for lower latency and load. Consequence: Hit rates improve cost, but correctness depends on keys, TTLs, and invalidation.
\par\textit{Caches trade staleness and invalidation complexity for lower latency and load. Tradeoff: Hit rates improve cost, but correctness depends on keys, TTLs, and invalidation. Watch for: Omitting tenant, model, prompt version, or authorization from a cache key leaks data.}
\paragraph{FC-1494 [medium]} Define Caching in interview-ready language.
\noindent\textbf{Answer.} Caches trade staleness and invalidation complexity for lower latency and load.
\par\textit{Caches trade staleness and invalidation complexity for lower latency and load.}
\paragraph{FC-1495 [hard]} State the main tradeoff and production pitfall for Caching.
\noindent\textbf{Answer.} Tradeoff: Hit rates improve cost, but correctness depends on keys, TTLs, and invalidation. Pitfall: Omitting tenant, model, prompt version, or authorization from a cache key leaks data.
\par\textit{Tradeoff: Hit rates improve cost, but correctness depends on keys, TTLs, and invalidation. Pitfall: Omitting tenant, model, prompt version, or authorization from a cache key leaks data.}
\paragraph{LA-1496 [hard]} Explain Caching from first principles, then show how you would evaluate and operate it in production.
\noindent\textbf{Answer.} Start with the mechanism: Caches trade staleness and invalidation complexity for lower latency and load. Make the tradeoff explicit: Hit rates improve cost, but correctness depends on keys, TTLs, and invalidation. Define workload-specific offline metrics and a latency/cost budget, test important slices, instrument the critical path, and create a rollback criterion. The production review must explicitly guard against this failure: Omitting tenant, model, prompt version, or authorization from a cache key leaks data.
\par\textit{A strong answer connects mechanism, measurable tradeoffs, failure isolation, observability, and rollback rather than listing product features.}
\paragraph{MCQ-1497 [medium]} A design requires this behavior: Balancers route by health, load, locality, affinity, or capacity across replicas. Which mechanism should the team study?
\begin{enumerate}[label=\Alph*.]
\item Exactly-once effects
\item Fault isolation
\item Load balancing
\item Multi-region design
\end{enumerate}
\noindent\textbf{Answer.} Load balancing
\par\textit{Balancers route by health, load, locality, affinity, or capacity across replicas. Tradeoff: Distribution improves availability but generic least-connections may ignore token workload size. Watch for: Sending a long-context request to the least-busy replica can still cause head-of-line blocking.}
\paragraph{MCQ-1498 [hard]} The team proposes Load balancing for a real workload. Which finding gives the correct decision boundary?
\begin{enumerate}[label=\Alph*.]
\item Business effects can be effectively once, but state and retention are needed.
\item Distribution improves availability but generic least-connections may ignore token workload size.
\item Blast radius shrinks at some efficiency cost.
\item Active-active improves availability and latency but makes writes and conflict resolution harder.
\end{enumerate}
\noindent\textbf{Answer.} Distribution improves availability but generic least-connections may ignore token workload size.
\par\textit{Balancers route by health, load, locality, affinity, or capacity across replicas. Tradeoff: Distribution improves availability but generic least-connections may ignore token workload size. Watch for: Sending a long-context request to the least-busy replica can still cause head-of-line blocking.}
\paragraph{MCQ-1499 [hard]} Before releasing Load balancing, which failure deserves a dedicated regression test?
\begin{enumerate}[label=\Alph*.]
\item A standby region that never receives realistic traffic may fail during disaster recovery.
\item Sending a long-context request to the least-busy replica can still cause head-of-line blocking.
\item Equating broker delivery semantics with end-to-end exactly-once effects is incorrect.
\item A single global vector index or queue can defeat otherwise isolated services.
\end{enumerate}
\noindent\textbf{Answer.} Sending a long-context request to the least-busy replica can still cause head-of-line blocking.
\par\textit{Balancers route by health, load, locality, affinity, or capacity across replicas. Tradeoff: Distribution improves availability but generic least-connections may ignore token workload size. Watch for: Sending a long-context request to the least-busy replica can still cause head-of-line blocking.}
\paragraph{MCQ-1500 [medium]} A design review reveals this mistake: Sending a long-context request to the least-busy replica can still cause head-of-line blocking. Which mechanism should be traced first?
\begin{enumerate}[label=\Alph*.]
\item Load balancing
\item Multi-region design
\item Fault isolation
\item Exactly-once effects
\end{enumerate}
\noindent\textbf{Answer.} Load balancing
\par\textit{Balancers route by health, load, locality, affinity, or capacity across replicas. Tradeoff: Distribution improves availability but generic least-connections may ignore token workload size. Watch for: Sending a long-context request to the least-busy replica can still cause head-of-line blocking.}
\paragraph{MCQ-1501 [hard]} Which causal explanation of Load balancing connects what it does to when it is worth its cost?
\begin{enumerate}[label=\Alph*.]
\item Mechanism: Systems typically combine at-least-once delivery with idempotent consumers, dedupe keys, and transactional outboxes. Consequence: Business effects can be effectively once, but state and retention are needed.
\item Mechanism: Regions trade user latency, residency, failover speed, replication lag, and operational complexity. Consequence: Active-active improves availability and latency but makes writes and conflict resolution harder.
\item Mechanism: Balancers route by health, load, locality, affinity, or capacity across replicas. Consequence: Distribution improves availability but generic least-connections may ignore token workload size.
\item Mechanism: Bulkheads, cells, quotas, and tenant boundaries prevent one failure or workload from consuming shared resources. Consequence: Blast radius shrinks at some efficiency cost.
\end{enumerate}
\noindent\textbf{Answer.} Mechanism: Balancers route by health, load, locality, affinity, or capacity across replicas. Consequence: Distribution improves availability but generic least-connections may ignore token workload size.
\par\textit{Balancers route by health, load, locality, affinity, or capacity across replicas. Tradeoff: Distribution improves availability but generic least-connections may ignore token workload size. Watch for: Sending a long-context request to the least-busy replica can still cause head-of-line blocking.}
\paragraph{FC-1502 [medium]} Define Load balancing in interview-ready language.
\noindent\textbf{Answer.} Balancers route by health, load, locality, affinity, or capacity across replicas.
\par\textit{Balancers route by health, load, locality, affinity, or capacity across replicas.}
\paragraph{FC-1503 [hard]} State the main tradeoff and production pitfall for Load balancing.
\noindent\textbf{Answer.} Tradeoff: Distribution improves availability but generic least-connections may ignore token workload size. Pitfall: Sending a long-context request to the least-busy replica can still cause head-of-line blocking.
\par\textit{Tradeoff: Distribution improves availability but generic least-connections may ignore token workload size. Pitfall: Sending a long-context request to the least-busy replica can still cause head-of-line blocking.}
\paragraph{LA-1504 [hard]} Explain Load balancing from first principles, then show how you would evaluate and operate it in production.
\noindent\textbf{Answer.} Start with the mechanism: Balancers route by health, load, locality, affinity, or capacity across replicas. Make the tradeoff explicit: Distribution improves availability but generic least-connections may ignore token workload size. Define workload-specific offline metrics and a latency/cost budget, test important slices, instrument the critical path, and create a rollback criterion. The production review must explicitly guard against this failure: Sending a long-context request to the least-busy replica can still cause head-of-line blocking.
\par\textit{A strong answer connects mechanism, measurable tradeoffs, failure isolation, observability, and rollback rather than listing product features.}
\paragraph{MCQ-1505 [medium]} A design requires this behavior: Bulkheads, cells, quotas, and tenant boundaries prevent one failure or workload from consuming shared resources. Which mechanism should the team study?
\begin{enumerate}[label=\Alph*.]
\item Retries and timeouts
\item Fault isolation
\item Exactly-once effects
\item Load balancing
\end{enumerate}
\noindent\textbf{Answer.} Fault isolation
\par\textit{Bulkheads, cells, quotas, and tenant boundaries prevent one failure or workload from consuming shared resources. Tradeoff: Blast radius shrinks at some efficiency cost. Watch for: A single global vector index or queue can defeat otherwise isolated services.}
\paragraph{MCQ-1506 [hard]} The team proposes Fault isolation for a real workload. Which finding gives the correct decision boundary?
\begin{enumerate}[label=\Alph*.]
\item Success rate improves, while retries amplify overload and tail latency.
\item Blast radius shrinks at some efficiency cost.
\item Business effects can be effectively once, but state and retention are needed.
\item Distribution improves availability but generic least-connections may ignore token workload size.
\end{enumerate}
\noindent\textbf{Answer.} Blast radius shrinks at some efficiency cost.
\par\textit{Bulkheads, cells, quotas, and tenant boundaries prevent one failure or workload from consuming shared resources. Tradeoff: Blast radius shrinks at some efficiency cost. Watch for: A single global vector index or queue can defeat otherwise isolated services.}
\paragraph{MCQ-1507 [hard]} Before releasing Fault isolation, which failure deserves a dedicated regression test?
\begin{enumerate}[label=\Alph*.]
\item A single global vector index or queue can defeat otherwise isolated services.
\item Sending a long-context request to the least-busy replica can still cause head-of-line blocking.
\item Equating broker delivery semantics with end-to-end exactly-once effects is incorrect.
\item Layered automatic retries create multiplicative request storms.
\end{enumerate}
\noindent\textbf{Answer.} A single global vector index or queue can defeat otherwise isolated services.
\par\textit{Bulkheads, cells, quotas, and tenant boundaries prevent one failure or workload from consuming shared resources. Tradeoff: Blast radius shrinks at some efficiency cost. Watch for: A single global vector index or queue can defeat otherwise isolated services.}
\paragraph{MCQ-1508 [medium]} A design review reveals this mistake: A single global vector index or queue can defeat otherwise isolated services. Which mechanism should be traced first?
\begin{enumerate}[label=\Alph*.]
\item Load balancing
\item Fault isolation
\item Retries and timeouts
\item Exactly-once effects
\end{enumerate}
\noindent\textbf{Answer.} Fault isolation
\par\textit{Bulkheads, cells, quotas, and tenant boundaries prevent one failure or workload from consuming shared resources. Tradeoff: Blast radius shrinks at some efficiency cost. Watch for: A single global vector index or queue can defeat otherwise isolated services.}
\paragraph{MCQ-1509 [hard]} Which causal explanation of Fault isolation connects what it does to when it is worth its cost?
\begin{enumerate}[label=\Alph*.]
\item Mechanism: Systems typically combine at-least-once delivery with idempotent consumers, dedupe keys, and transactional outboxes. Consequence: Business effects can be effectively once, but state and retention are needed.
\item Mechanism: Balancers route by health, load, locality, affinity, or capacity across replicas. Consequence: Distribution improves availability but generic least-connections may ignore token workload size.
\item Mechanism: Timeouts bound uncertainty; retries use budgets, backoff, jitter, and idempotency for transient failures. Consequence: Success rate improves, while retries amplify overload and tail latency.
\item Mechanism: Bulkheads, cells, quotas, and tenant boundaries prevent one failure or workload from consuming shared resources. Consequence: Blast radius shrinks at some efficiency cost.
\end{enumerate}
\noindent\textbf{Answer.} Mechanism: Bulkheads, cells, quotas, and tenant boundaries prevent one failure or workload from consuming shared resources. Consequence: Blast radius shrinks at some efficiency cost.
\par\textit{Bulkheads, cells, quotas, and tenant boundaries prevent one failure or workload from consuming shared resources. Tradeoff: Blast radius shrinks at some efficiency cost. Watch for: A single global vector index or queue can defeat otherwise isolated services.}
\paragraph{FC-1510 [medium]} Define Fault isolation in interview-ready language.
\noindent\textbf{Answer.} Bulkheads, cells, quotas, and tenant boundaries prevent one failure or workload from consuming shared resources.
\par\textit{Bulkheads, cells, quotas, and tenant boundaries prevent one failure or workload from consuming shared resources.}
\paragraph{FC-1511 [hard]} State the main tradeoff and production pitfall for Fault isolation.
\noindent\textbf{Answer.} Tradeoff: Blast radius shrinks at some efficiency cost. Pitfall: A single global vector index or queue can defeat otherwise isolated services.
\par\textit{Tradeoff: Blast radius shrinks at some efficiency cost. Pitfall: A single global vector index or queue can defeat otherwise isolated services.}
\paragraph{LA-1512 [hard]} Explain Fault isolation from first principles, then show how you would evaluate and operate it in production.
\noindent\textbf{Answer.} Start with the mechanism: Bulkheads, cells, quotas, and tenant boundaries prevent one failure or workload from consuming shared resources. Make the tradeoff explicit: Blast radius shrinks at some efficiency cost. Define workload-specific offline metrics and a latency/cost budget, test important slices, instrument the critical path, and create a rollback criterion. The production review must explicitly guard against this failure: A single global vector index or queue can defeat otherwise isolated services.
\par\textit{A strong answer connects mechanism, measurable tradeoffs, failure isolation, observability, and rollback rather than listing product features.}
\paragraph{MCQ-1513 [medium]} A design requires this behavior: Systems typically combine at-least-once delivery with idempotent consumers, dedupe keys, and transactional outboxes. Which mechanism should the team study?
\begin{enumerate}[label=\Alph*.]
\item CAP and consistency
\item Retries and timeouts
\item Exactly-once effects
\item Caching
\end{enumerate}
\noindent\textbf{Answer.} Exactly-once effects
\par\textit{Systems typically combine at-least-once delivery with idempotent consumers, dedupe keys, and transactional outboxes. Tradeoff: Business effects can be effectively once, but state and retention are needed. Watch for: Equating broker delivery semantics with end-to-end exactly-once effects is incorrect.}
\paragraph{MCQ-1514 [hard]} The team proposes Exactly-once effects for a real workload. Which finding gives the correct decision boundary?
\begin{enumerate}[label=\Alph*.]
\item Stronger guarantees simplify application invariants but increase latency and unavailability.
\item Hit rates improve cost, but correctness depends on keys, TTLs, and invalidation.
\item Business effects can be effectively once, but state and retention are needed.
\item Success rate improves, while retries amplify overload and tail latency.
\end{enumerate}
\noindent\textbf{Answer.} Business effects can be effectively once, but state and retention are needed.
\par\textit{Systems typically combine at-least-once delivery with idempotent consumers, dedupe keys, and transactional outboxes. Tradeoff: Business effects can be effectively once, but state and retention are needed. Watch for: Equating broker delivery semantics with end-to-end exactly-once effects is incorrect.}
\paragraph{MCQ-1515 [hard]} Before releasing Exactly-once effects, which failure deserves a dedicated regression test?
\begin{enumerate}[label=\Alph*.]
\item Quoting CAP without identifying partitions and operations does not guide design.
\item Equating broker delivery semantics with end-to-end exactly-once effects is incorrect.
\item Omitting tenant, model, prompt version, or authorization from a cache key leaks data.
\item Layered automatic retries create multiplicative request storms.
\end{enumerate}
\noindent\textbf{Answer.} Equating broker delivery semantics with end-to-end exactly-once effects is incorrect.
\par\textit{Systems typically combine at-least-once delivery with idempotent consumers, dedupe keys, and transactional outboxes. Tradeoff: Business effects can be effectively once, but state and retention are needed. Watch for: Equating broker delivery semantics with end-to-end exactly-once effects is incorrect.}
\paragraph{MCQ-1516 [medium]} A design review reveals this mistake: Equating broker delivery semantics with end-to-end exactly-once effects is incorrect. Which mechanism should be traced first?
\begin{enumerate}[label=\Alph*.]
\item Caching
\item CAP and consistency
\item Retries and timeouts
\item Exactly-once effects
\end{enumerate}
\noindent\textbf{Answer.} Exactly-once effects
\par\textit{Systems typically combine at-least-once delivery with idempotent consumers, dedupe keys, and transactional outboxes. Tradeoff: Business effects can be effectively once, but state and retention are needed. Watch for: Equating broker delivery semantics with end-to-end exactly-once effects is incorrect.}
\paragraph{MCQ-1517 [hard]} Which causal explanation of Exactly-once effects connects what it does to when it is worth its cost?
\begin{enumerate}[label=\Alph*.]
\item Mechanism: Timeouts bound uncertainty; retries use budgets, backoff, jitter, and idempotency for transient failures. Consequence: Success rate improves, while retries amplify overload and tail latency.
\item Mechanism: Caches trade staleness and invalidation complexity for lower latency and load. Consequence: Hit rates improve cost, but correctness depends on keys, TTLs, and invalidation.
\item Mechanism: Under a network partition, a distributed system must trade availability against consistent reads and writes. Consequence: Stronger guarantees simplify application invariants but increase latency and unavailability.
\item Mechanism: Systems typically combine at-least-once delivery with idempotent consumers, dedupe keys, and transactional outboxes. Consequence: Business effects can be effectively once, but state and retention are needed.
\end{enumerate}
\noindent\textbf{Answer.} Mechanism: Systems typically combine at-least-once delivery with idempotent consumers, dedupe keys, and transactional outboxes. Consequence: Business effects can be effectively once, but state and retention are needed.
\par\textit{Systems typically combine at-least-once delivery with idempotent consumers, dedupe keys, and transactional outboxes. Tradeoff: Business effects can be effectively once, but state and retention are needed. Watch for: Equating broker delivery semantics with end-to-end exactly-once effects is incorrect.}
\paragraph{FC-1518 [medium]} Define Exactly-once effects in interview-ready language.
\noindent\textbf{Answer.} Systems typically combine at-least-once delivery with idempotent consumers, dedupe keys, and transactional outboxes.
\par\textit{Systems typically combine at-least-once delivery with idempotent consumers, dedupe keys, and transactional outboxes.}
\paragraph{FC-1519 [hard]} State the main tradeoff and production pitfall for Exactly-once effects.
\noindent\textbf{Answer.} Tradeoff: Business effects can be effectively once, but state and retention are needed. Pitfall: Equating broker delivery semantics with end-to-end exactly-once effects is incorrect.
\par\textit{Tradeoff: Business effects can be effectively once, but state and retention are needed. Pitfall: Equating broker delivery semantics with end-to-end exactly-once effects is incorrect.}
\paragraph{LA-1520 [hard]} Explain Exactly-once effects from first principles, then show how you would evaluate and operate it in production.
\noindent\textbf{Answer.} Start with the mechanism: Systems typically combine at-least-once delivery with idempotent consumers, dedupe keys, and transactional outboxes. Make the tradeoff explicit: Business effects can be effectively once, but state and retention are needed. Define workload-specific offline metrics and a latency/cost budget, test important slices, instrument the critical path, and create a rollback criterion. The production review must explicitly guard against this failure: Equating broker delivery semantics with end-to-end exactly-once effects is incorrect.
\par\textit{A strong answer connects mechanism, measurable tradeoffs, failure isolation, observability, and rollback rather than listing product features.}
\paragraph{MCQ-1521 [medium]} A design requires this behavior: Regions trade user latency, residency, failover speed, replication lag, and operational complexity. Which mechanism should the team study?
\begin{enumerate}[label=\Alph*.]
\item Consensus
\item Multi-region design
\item Load balancing
\item Exactly-once effects
\end{enumerate}
\noindent\textbf{Answer.} Multi-region design
\par\textit{Regions trade user latency, residency, failover speed, replication lag, and operational complexity. Tradeoff: Active-active improves availability and latency but makes writes and conflict resolution harder. Watch for: A standby region that never receives realistic traffic may fail during disaster recovery.}
\paragraph{MCQ-1522 [hard]} The team proposes Multi-region design for a real workload. Which finding gives the correct decision boundary?
\begin{enumerate}[label=\Alph*.]
\item Active-active improves availability and latency but makes writes and conflict resolution harder.
\item Business effects can be effectively once, but state and retention are needed.
\item Coordination provides a consistent control plane but costs extra messages and quorum latency.
\item Distribution improves availability but generic least-connections may ignore token workload size.
\end{enumerate}
\noindent\textbf{Answer.} Active-active improves availability and latency but makes writes and conflict resolution harder.
\par\textit{Regions trade user latency, residency, failover speed, replication lag, and operational complexity. Tradeoff: Active-active improves availability and latency but makes writes and conflict resolution harder. Watch for: A standby region that never receives realistic traffic may fail during disaster recovery.}
\paragraph{MCQ-1523 [hard]} Before releasing Multi-region design, which failure deserves a dedicated regression test?
\begin{enumerate}[label=\Alph*.]
\item Sending a long-context request to the least-busy replica can still cause head-of-line blocking.
\item Equating broker delivery semantics with end-to-end exactly-once effects is incorrect.
\item A standby region that never receives realistic traffic may fail during disaster recovery.
\item Deploying a quorum across unreliable links can reduce availability.
\end{enumerate}
\noindent\textbf{Answer.} A standby region that never receives realistic traffic may fail during disaster recovery.
\par\textit{Regions trade user latency, residency, failover speed, replication lag, and operational complexity. Tradeoff: Active-active improves availability and latency but makes writes and conflict resolution harder. Watch for: A standby region that never receives realistic traffic may fail during disaster recovery.}
\paragraph{MCQ-1524 [medium]} A design review reveals this mistake: A standby region that never receives realistic traffic may fail during disaster recovery. Which mechanism should be traced first?
\begin{enumerate}[label=\Alph*.]
\item Multi-region design
\item Exactly-once effects
\item Load balancing
\item Consensus
\end{enumerate}
\noindent\textbf{Answer.} Multi-region design
\par\textit{Regions trade user latency, residency, failover speed, replication lag, and operational complexity. Tradeoff: Active-active improves availability and latency but makes writes and conflict resolution harder. Watch for: A standby region that never receives realistic traffic may fail during disaster recovery.}
\paragraph{MCQ-1525 [hard]} Which causal explanation of Multi-region design connects what it does to when it is worth its cost?
\begin{enumerate}[label=\Alph*.]
\item Mechanism: Regions trade user latency, residency, failover speed, replication lag, and operational complexity. Consequence: Active-active improves availability and latency but makes writes and conflict resolution harder.
\item Mechanism: Protocols such as Raft replicate an ordered log so nodes agree despite failures. Consequence: Coordination provides a consistent control plane but costs extra messages and quorum latency.
\item Mechanism: Balancers route by health, load, locality, affinity, or capacity across replicas. Consequence: Distribution improves availability but generic least-connections may ignore token workload size.
\item Mechanism: Systems typically combine at-least-once delivery with idempotent consumers, dedupe keys, and transactional outboxes. Consequence: Business effects can be effectively once, but state and retention are needed.
\end{enumerate}
\noindent\textbf{Answer.} Mechanism: Regions trade user latency, residency, failover speed, replication lag, and operational complexity. Consequence: Active-active improves availability and latency but makes writes and conflict resolution harder.
\par\textit{Regions trade user latency, residency, failover speed, replication lag, and operational complexity. Tradeoff: Active-active improves availability and latency but makes writes and conflict resolution harder. Watch for: A standby region that never receives realistic traffic may fail during disaster recovery.}
\paragraph{FC-1526 [medium]} Define Multi-region design in interview-ready language.
\noindent\textbf{Answer.} Regions trade user latency, residency, failover speed, replication lag, and operational complexity.
\par\textit{Regions trade user latency, residency, failover speed, replication lag, and operational complexity.}
\paragraph{FC-1527 [hard]} State the main tradeoff and production pitfall for Multi-region design.
\noindent\textbf{Answer.} Tradeoff: Active-active improves availability and latency but makes writes and conflict resolution harder. Pitfall: A standby region that never receives realistic traffic may fail during disaster recovery.
\par\textit{Tradeoff: Active-active improves availability and latency but makes writes and conflict resolution harder. Pitfall: A standby region that never receives realistic traffic may fail during disaster recovery.}
\paragraph{LA-1528 [hard]} Explain Multi-region design from first principles, then show how you would evaluate and operate it in production.
\noindent\textbf{Answer.} Start with the mechanism: Regions trade user latency, residency, failover speed, replication lag, and operational complexity. Make the tradeoff explicit: Active-active improves availability and latency but makes writes and conflict resolution harder. Define workload-specific offline metrics and a latency/cost budget, test important slices, instrument the critical path, and create a rollback criterion. The production review must explicitly guard against this failure: A standby region that never receives realistic traffic may fail during disaster recovery.
\par\textit{A strong answer connects mechanism, measurable tradeoffs, failure isolation, observability, and rollback rather than listing product features.}
\paragraph{MCQ-1529 [medium]} A design requires this behavior: Workload models connect arrival rate, prompt and output tokens, service time, concurrency, memory, and SLO targets. Which mechanism should the team study?
\begin{enumerate}[label=\Alph*.]
\item Capacity planning
\item Retries and timeouts
\item Multi-region design
\item Queues and backpressure
\end{enumerate}
\noindent\textbf{Answer.} Capacity planning
\par\textit{Workload models connect arrival rate, prompt and output tokens, service time, concurrency, memory, and SLO targets. Tradeoff: Headroom handles bursts but idle accelerators are expensive. Watch for: Planning with average sequence length underestimates tail memory and latency.}
\paragraph{MCQ-1530 [hard]} The team proposes Capacity planning for a real workload. Which finding gives the correct decision boundary?
\begin{enumerate}[label=\Alph*.]
\item Active-active improves availability and latency but makes writes and conflict resolution harder.
\item Smoothing improves utilization, while waiting increases latency and stale work.
\item Headroom handles bursts but idle accelerators are expensive.
\item Success rate improves, while retries amplify overload and tail latency.
\end{enumerate}
\noindent\textbf{Answer.} Headroom handles bursts but idle accelerators are expensive.
\par\textit{Workload models connect arrival rate, prompt and output tokens, service time, concurrency, memory, and SLO targets. Tradeoff: Headroom handles bursts but idle accelerators are expensive. Watch for: Planning with average sequence length underestimates tail memory and latency.}
\paragraph{MCQ-1531 [hard]} Before releasing Capacity planning, which failure deserves a dedicated regression test?
\begin{enumerate}[label=\Alph*.]
\item Layered automatic retries create multiplicative request storms.
\item Planning with average sequence length underestimates tail memory and latency.
\item An unbounded queue converts overload into memory exhaustion.
\item A standby region that never receives realistic traffic may fail during disaster recovery.
\end{enumerate}
\noindent\textbf{Answer.} Planning with average sequence length underestimates tail memory and latency.
\par\textit{Workload models connect arrival rate, prompt and output tokens, service time, concurrency, memory, and SLO targets. Tradeoff: Headroom handles bursts but idle accelerators are expensive. Watch for: Planning with average sequence length underestimates tail memory and latency.}
\paragraph{MCQ-1532 [medium]} A design review reveals this mistake: Planning with average sequence length underestimates tail memory and latency. Which mechanism should be traced first?
\begin{enumerate}[label=\Alph*.]
\item Capacity planning
\item Retries and timeouts
\item Queues and backpressure
\item Multi-region design
\end{enumerate}
\noindent\textbf{Answer.} Capacity planning
\par\textit{Workload models connect arrival rate, prompt and output tokens, service time, concurrency, memory, and SLO targets. Tradeoff: Headroom handles bursts but idle accelerators are expensive. Watch for: Planning with average sequence length underestimates tail memory and latency.}
\paragraph{MCQ-1533 [hard]} Which causal explanation of Capacity planning connects what it does to when it is worth its cost?
\begin{enumerate}[label=\Alph*.]
\item Mechanism: Timeouts bound uncertainty; retries use budgets, backoff, jitter, and idempotency for transient failures. Consequence: Success rate improves, while retries amplify overload and tail latency.
\item Mechanism: Workload models connect arrival rate, prompt and output tokens, service time, concurrency, memory, and SLO targets. Consequence: Headroom handles bursts but idle accelerators are expensive.
\item Mechanism: Bounded queues absorb bursts while admission control and backpressure prevent producers from overwhelming consumers. Consequence: Smoothing improves utilization, while waiting increases latency and stale work.
\item Mechanism: Regions trade user latency, residency, failover speed, replication lag, and operational complexity. Consequence: Active-active improves availability and latency but makes writes and conflict resolution harder.
\end{enumerate}
\noindent\textbf{Answer.} Mechanism: Workload models connect arrival rate, prompt and output tokens, service time, concurrency, memory, and SLO targets. Consequence: Headroom handles bursts but idle accelerators are expensive.
\par\textit{Workload models connect arrival rate, prompt and output tokens, service time, concurrency, memory, and SLO targets. Tradeoff: Headroom handles bursts but idle accelerators are expensive. Watch for: Planning with average sequence length underestimates tail memory and latency.}
\paragraph{FC-1534 [medium]} Define Capacity planning in interview-ready language.
\noindent\textbf{Answer.} Workload models connect arrival rate, prompt and output tokens, service time, concurrency, memory, and SLO targets.
\par\textit{Workload models connect arrival rate, prompt and output tokens, service time, concurrency, memory, and SLO targets.}
\paragraph{FC-1535 [hard]} State the main tradeoff and production pitfall for Capacity planning.
\noindent\textbf{Answer.} Tradeoff: Headroom handles bursts but idle accelerators are expensive. Pitfall: Planning with average sequence length underestimates tail memory and latency.
\par\textit{Tradeoff: Headroom handles bursts but idle accelerators are expensive. Pitfall: Planning with average sequence length underestimates tail memory and latency.}
\paragraph{LA-1536 [hard]} Explain Capacity planning from first principles, then show how you would evaluate and operate it in production.
\noindent\textbf{Answer.} Start with the mechanism: Workload models connect arrival rate, prompt and output tokens, service time, concurrency, memory, and SLO targets. Make the tradeoff explicit: Headroom handles bursts but idle accelerators are expensive. Define workload-specific offline metrics and a latency/cost budget, test important slices, instrument the critical path, and create a rollback criterion. The production review must explicitly guard against this failure: Planning with average sequence length underestimates tail memory and latency.
\par\textit{A strong answer connects mechanism, measurable tradeoffs, failure isolation, observability, and rollback rather than listing product features.}
\clearpage\section{Primary sources}
\begin{enumerate}
\item \href{https://arxiv.org/abs/1706.03762}{Attention Is All You Need} (paper)
\item \href{https://arxiv.org/abs/2005.11401}{Retrieval-Augmented Generation for Knowledge-Intensive NLP Tasks} (paper)
\item \href{https://arxiv.org/abs/2004.04906}{Dense Passage Retrieval for Open-Domain Question Answering} (paper)
\item \href{https://arxiv.org/abs/2104.08663}{BEIR: A Heterogeneous Benchmark for Zero-shot Evaluation of Information Retrieval Models} (paper)
\item \href{https://arxiv.org/abs/2309.15217}{RAGAS: Automated Evaluation of Retrieval Augmented Generation} (paper)
\item \href{https://arxiv.org/abs/2109.07958}{TruthfulQA: Measuring How Models Mimic Human Falsehoods} (paper)
\item \href{https://arxiv.org/abs/2210.03629}{ReAct: Synergizing Reasoning and Acting in Language Models} (paper)
\item \href{https://arxiv.org/abs/2303.11366}{Reflexion: Language Agents with Verbal Reinforcement Learning} (paper)
\item \href{https://arxiv.org/abs/2302.04761}{Toolformer: Language Models Can Teach Themselves to Use Tools} (paper)
\item \href{https://arxiv.org/abs/2404.16130}{A Graph RAG Approach to Query-Focused Summarization} (paper)
\item \href{https://arxiv.org/abs/2409.13731}{KAG: Boosting LLMs in Professional Domains via Knowledge Augmented Generation} (paper)
\item \href{https://modelcontextprotocol.io/specification/2026-07-28}{Model Context Protocol Specification (2026-07-28)} (specification)
\item \href{https://arxiv.org/abs/2205.14135}{FlashAttention: Fast and Memory-Efficient Exact Attention with IO-Awareness} (paper)
\item \href{https://arxiv.org/abs/2309.06180}{Efficient Memory Management for Large Language Model Serving with PagedAttention} (paper)
\item \href{https://arxiv.org/abs/2305.13245}{GQA: Training Generalized Multi-Query Transformer Models} (paper)
\item \href{https://arxiv.org/abs/2211.17192}{Fast Inference from Transformers via Speculative Decoding} (paper)
\item \href{https://arxiv.org/abs/2106.09685}{LoRA: Low-Rank Adaptation of Large Language Models} (paper)
\item \href{https://arxiv.org/abs/2305.14314}{QLoRA: Efficient Finetuning of Quantized LLMs} (paper)
\item \href{https://arxiv.org/abs/2305.18290}{Direct Preference Optimization} (paper)
\item \href{https://github.com/pgvector/pgvector}{pgvector Documentation} (official docs)
\item \href{https://docs.pinecone.io/guides/search/hybrid-search}{Pinecone Hybrid Search} (official docs)
\item \href{https://qdrant.tech/documentation/}{Qdrant Documentation} (official docs)
\item \href{https://docs.weaviate.io/weaviate}{Weaviate Documentation} (official docs)
\item \href{https://milvus.io/docs}{Milvus Documentation} (official docs)
\item \href{https://www.elastic.co/docs/solutions/search/vector/knn}{Elasticsearch kNN Search} (official docs)
\item \href{https://neo4j.com/docs/cypher-manual/current/indexes/semantic-indexes/vector-indexes/}{Neo4j Vector Indexes} (official docs)
\item \href{https://docs.langchain.com/oss/python/langgraph/overview}{LangGraph Overview} (official docs)
\item \href{https://docs.temporal.io/}{Temporal Durable Execution} (official docs)
\item \href{https://opentelemetry.io/docs/specs/semconv/}{OpenTelemetry Semantic Conventions} (standard)
\item \href{https://mlflow.org/docs/latest/genai/}{MLflow for LLMs and Agents} (official docs)
\item \href{https://docs.evidentlyai.com/}{Evidently Documentation} (official docs)
\item \href{https://doc.dvc.org/start/data-pipelines/data-pipelines}{DVC Data Pipelines and Versioning} (official docs)
\item \href{https://docs.feast.dev/}{Feast Feature Store Documentation} (official docs)
\item \href{https://docs.vllm.ai/}{vLLM Documentation} (official docs)
\item \href{https://docs.ray.io/en/latest/serve/llm/}{Ray Serve LLM Serving} (official docs)
\item \href{https://kserve.github.io/website/}{KServe Documentation} (official docs)
\item \href{https://kubernetes.io/docs/tasks/run-application/horizontal-pod-autoscale/}{Kubernetes Horizontal Pod Autoscaling} (official docs)
\item \href{https://fastapi.tiangolo.com/}{FastAPI Documentation} (official docs)
\item \href{https://grpc.io/docs/}{gRPC Documentation} (official docs)
\item \href{https://docs.python.org/3/}{Python 3 Documentation} (official docs)
\item \href{https://genai.owasp.org/llm-top-10/}{OWASP Top 10 for LLM Applications} (standard)
\item \href{https://www.nist.gov/itl/ai-risk-management-framework}{NIST AI Risk Management Framework} (standard)
\item \href{https://docs.langchain.com/oss/python/langchain/overview}{LangChain official documentation} (official technology documentation)
\item \href{https://docs.langchain.com/oss/python/langchain/agents}{LangChain agents and runtime} (official technology documentation)
\item \href{https://docs.langchain.com/oss/python/langgraph/graph-api}{LangGraph graph API} (official technology documentation)
\item \href{https://docs.langchain.com/oss/python/langgraph/persistence}{LangGraph persistence design} (official technology documentation)
\item \href{https://docs.langchain.com/oss/python/deepagents/overview}{Deep Agents official documentation} (official technology documentation)
\item \href{https://docs.langchain.com/oss/python/deepagents/harness}{Deep Agents harness} (official technology documentation)
\item \href{https://docs.langchain.com/oss/python/deepagents/backends}{Deep Agents backends} (official technology documentation)
\item \href{https://docs.llamaindex.ai/en/stable/}{LlamaIndex official documentation} (official technology documentation)
\item \href{https://docs.haystack.deepset.ai/docs/intro}{Haystack official documentation} (official technology documentation)
\item \href{https://learn.microsoft.com/en-us/semantic-kernel/overview/}{Semantic Kernel official documentation} (official technology documentation)
\item \href{https://microsoft.github.io/autogen/stable/}{AutoGen official documentation} (official technology documentation)
\item \href{https://docs.crewai.com/}{CrewAI official documentation} (official technology documentation)
\item \href{https://ai.pydantic.dev/}{PydanticAI official documentation} (official technology documentation)
\item \href{https://dspy.ai/}{DSPy official documentation} (official technology documentation)
\item \href{https://openai.github.io/openai-agents-python/}{OpenAI Agents SDK official documentation} (official technology documentation)
\item \href{https://google.github.io/adk-docs/}{Google Agent Development Kit official documentation} (official technology documentation)
\item \href{https://mastra.ai/docs}{Mastra official documentation} (official technology documentation)
\item \href{https://docs.vllm.ai/en/latest/design/arch_overview/}{vLLM architecture} (official technology documentation)
\item \href{https://docs.vllm.ai/en/latest/serving/openai_compatible_server/}{vLLM serving guide} (official technology documentation)
\item \href{https://docs.sglang.ai/}{SGLang official documentation} (official technology documentation)
\item \href{https://huggingface.co/docs/text-generation-inference/index}{Hugging Face Text Generation Inference official documentation} (official technology documentation)
\item \href{https://nvidia.github.io/TensorRT-LLM/}{TensorRT-LLM official documentation} (official technology documentation)
\item \href{https://docs.nvidia.com/deeplearning/triton-inference-server/user-guide/docs/}{NVIDIA Triton Inference Server official documentation} (official technology documentation)
\item \href{https://github.com/ggml-org/llama.cpp}{llama.cpp official documentation} (official technology documentation)
\item \href{https://docs.ollama.com/}{Ollama official documentation} (official technology documentation)
\item \href{https://docs.ray.io/en/latest/serve/}{Ray Serve official documentation} (official technology documentation)
\item \href{https://docs.bentoml.com/}{BentoML official documentation} (official technology documentation)
\item \href{https://localai.io/}{LocalAI official documentation} (official technology documentation)
\item \href{https://modal.com/docs/guide}{Modal official documentation} (official technology documentation)
\item \href{https://docs.baseten.co/}{Baseten official documentation} (official technology documentation)
\item \href{https://huggingface.co/docs/inference-endpoints/index}{Hugging Face Inference Endpoints official documentation} (official technology documentation)
\item \href{https://docs.pinecone.io/}{Pinecone official documentation} (official technology documentation)
\item \href{https://docs.trychroma.com/}{Chroma official documentation} (official technology documentation)
\item \href{https://docs.lancedb.com/}{LanceDB official documentation} (official technology documentation)
\item \href{https://github.com/facebookresearch/faiss}{FAISS official documentation} (official technology documentation)
\item \href{https://www.elastic.co/docs/solutions/search}{Elasticsearch official documentation} (official technology documentation)
\item \href{https://docs.opensearch.org/latest/vector-search/}{OpenSearch official documentation} (official technology documentation)
\item \href{https://docs.vespa.ai/}{Vespa official documentation} (official technology documentation)
\item \href{https://redis.io/docs/latest/develop/ai/}{Redis Search and Vector Sets official documentation} (official technology documentation)
\item \href{https://neo4j.com/docs/}{Neo4j official documentation} (official technology documentation)
\item \href{https://neo4j.com/docs/getting-started/appendix/graphdb-concepts/}{Neo4j graph database concepts} (official technology documentation)
\item \href{https://neo4j.com/docs/cypher-manual/current/}{Cypher manual} (official technology documentation)
\item \href{https://neo4j.com/docs/neo4j-graphrag-python/current/}{Neo4j GraphRAG for Python} (official technology documentation)
\item \href{https://memgraph.com/docs}{Memgraph official documentation} (official technology documentation)
\item \href{https://docs.falkordb.com/}{FalkorDB official documentation} (official technology documentation)
\item \href{https://docs.aws.amazon.com/neptune/}{Amazon Neptune official documentation} (official technology documentation)
\item \href{https://docs.kuzudb.com/}{Kuzu official documentation} (official technology documentation)
\item \href{https://docs.langchain.com/langsmith/home}{LangSmith official documentation} (official technology documentation)
\item \href{https://docs.langchain.com/langsmith/observability}{LangSmith observability} (official technology documentation)
\item \href{https://docs.langchain.com/langsmith/evaluation}{LangSmith evaluation} (official technology documentation)
\item \href{https://langfuse.com/docs}{Langfuse official documentation} (official technology documentation)
\item \href{https://arize.com/docs/phoenix}{Arize Phoenix official documentation} (official technology documentation)
\item \href{https://docs.helicone.ai/}{Helicone official documentation} (official technology documentation)
\item \href{https://opentelemetry.io/docs/}{OpenTelemetry official documentation} (official technology documentation)
\item \href{https://weave-docs.wandb.ai/}{Weights \& Biases Weave official documentation} (official technology documentation)
\item \href{https://www.braintrust.dev/docs}{Braintrust official documentation} (official technology documentation)
\item \href{https://docs.ragas.io/}{Ragas official documentation} (official technology documentation)
\item \href{https://deepeval.com/docs/}{DeepEval official documentation} (official technology documentation)
\item \href{https://www.trulens.org/}{TruLens official documentation} (official technology documentation)
\item \href{https://www.promptfoo.dev/docs/}{Promptfoo official documentation} (official technology documentation)
\item \href{https://modelcontextprotocol.io/docs/learn/architecture}{MCP architecture} (official technology documentation)
\item \href{https://modelcontextprotocol.io/specification/latest/basic/authorization}{MCP authorization} (official technology documentation)
\item \href{https://modelcontextprotocol.io/docs/sdk}{Official MCP SDKs official documentation} (official technology documentation)
\item \href{https://gofastmcp.com/getting-started/welcome}{FastMCP official documentation} (official technology documentation)
\item \href{https://github.com/modelcontextprotocol/inspector}{MCP Inspector official documentation} (official technology documentation)
\item \href{https://docs.docker.com/ai/mcp-catalog-and-toolkit/}{Docker MCP Toolkit and Gateway official documentation} (official technology documentation)
\item \href{https://smithery.ai/docs}{Smithery official documentation} (official technology documentation)
\item \href{https://docs.composio.dev/}{Composio official documentation} (official technology documentation)
\item \href{https://docs.arcade.dev/}{Arcade official documentation} (official technology documentation)
\item \href{https://e2b.dev/docs}{E2B official documentation} (official technology documentation)
\item \href{https://www.daytona.io/docs/}{Daytona official documentation} (official technology documentation)
\item \href{https://modal.com/docs/guide/sandbox}{Modal Sandboxes official documentation} (official technology documentation)
\item \href{https://docs.docker.com/engine/security/}{Docker official documentation} (official technology documentation)
\item \href{https://gvisor.dev/docs/}{gVisor official documentation} (official technology documentation)
\item \href{https://github.com/firecracker-microvm/firecracker}{Firecracker official documentation} (official technology documentation)
\item \href{https://katacontainers.io/docs/}{Kata Containers official documentation} (official technology documentation)
\item \href{https://kubernetes.io/docs/concepts/workloads/controllers/job/}{Kubernetes Jobs official documentation} (official technology documentation)
\item \href{https://airflow.apache.org/docs/}{Apache Airflow official documentation} (official technology documentation)
\item \href{https://docs.prefect.io/}{Prefect official documentation} (official technology documentation)
\item \href{https://docs.dagster.io/}{Dagster official documentation} (official technology documentation)
\item \href{https://dvc.org/doc}{DVC official documentation} (official technology documentation)
\item \href{https://docs.lakefs.io/}{lakeFS official documentation} (official technology documentation)
\item \href{https://www.kubeflow.org/docs/components/pipelines/}{Kubeflow Pipelines official documentation} (official technology documentation)
\item \href{https://docs.zenml.io/}{ZenML official documentation} (official technology documentation)
\item \href{https://help.getzep.com/}{Zep official documentation} (official technology documentation)
\item \href{https://docs.mem0.ai/}{Mem0 official documentation} (official technology documentation)
\item \href{https://docs.pytorch.org/docs/stable/index.html}{PyTorch official documentation} (official technology documentation)
\item \href{https://docs.jax.dev/en/latest/}{JAX official documentation} (official technology documentation)
\item \href{https://huggingface.co/docs/transformers/index}{Hugging Face Transformers official documentation} (official technology documentation)
\item \href{https://sbert.net/}{Sentence Transformers official documentation} (official technology documentation)
\item \href{https://huggingface.co/docs/peft/index}{Hugging Face PEFT official documentation} (official technology documentation)
\item \href{https://huggingface.co/docs/bitsandbytes/main/en/index}{bitsandbytes official documentation} (official technology documentation)
\item \href{https://onnxruntime.ai/docs/}{ONNX Runtime official documentation} (official technology documentation)
\item \href{https://docs.openvino.ai/}{OpenVINO official documentation} (official technology documentation)
\item \href{https://bge-model.com/}{FlagEmbedding official documentation} (official technology documentation)
\item \href{https://python.useinstructor.com/}{Instructor official documentation} (official technology documentation)
\item \href{https://docs.cohere.com/docs/rerank-overview}{Cohere Rerank official documentation} (official technology documentation)
\item \href{https://huggingface.co/docs/text-embeddings-inference/index}{Hugging Face Text Embeddings Inference official documentation} (official technology documentation)
\item \href{https://docs.litellm.ai/}{LiteLLM official documentation} (official technology documentation)
\item \href{https://portkey.ai/docs/}{Portkey AI Gateway official documentation} (official technology documentation)
\item \href{https://developer.konghq.com/ai-gateway/}{Kong AI Gateway official documentation} (official technology documentation)
\item \href{https://aigateway.envoyproxy.io/docs/}{Envoy AI Gateway official documentation} (official technology documentation)
\item \href{https://docs.aws.amazon.com/sagemaker/latest/dg/deploy-model.html}{Amazon SageMaker AI Endpoints official documentation} (official technology documentation)
\item \href{https://cloud.google.com/vertex-ai/docs/predictions/overview}{Vertex AI Endpoints official documentation} (official technology documentation)
\item \href{https://learn.microsoft.com/en-us/azure/machine-learning/concept-endpoints-online}{Azure Machine Learning Managed Online Endpoints official documentation} (official technology documentation)
\item \href{https://docs.nvidia.com/nemo/guardrails/latest/}{NVIDIA NeMo Guardrails official documentation} (official technology documentation)
\item \href{https://www.guardrailsai.com/docs}{Guardrails AI official documentation} (official technology documentation)
\item \href{https://www.llama.com/docs/model-cards-and-prompt-formats/llama-guard-4/}{Llama Guard official documentation} (official technology documentation)
\item \href{https://protectai.github.io/llm-guard/}{Protect AI LLM Guard official documentation} (official technology documentation)
\item \href{https://microsoft.github.io/presidio/}{Microsoft Presidio official documentation} (official technology documentation)
\item \href{https://www.openpolicyagent.org/docs/latest/}{Open Policy Agent official documentation} (official technology documentation)
\item \href{https://docs.cedarpolicy.com/}{Cedar official documentation} (official technology documentation)
\item \href{https://docs.lakera.ai/docs}{Lakera Guard official documentation} (official technology documentation)
\item \href{https://docs.agno.com/}{Agno official documentation} (official technology documentation)
\item \href{https://huggingface.co/docs/smolagents/index}{Hugging Face smolagents official documentation} (official technology documentation)
\item \href{https://docs.letta.com/}{Letta official documentation} (official technology documentation)
\item \href{https://langroid.github.io/langroid/}{Langroid official documentation} (official technology documentation)
\item \href{https://framework.beeai.dev/}{BeeAI Framework official documentation} (official technology documentation)
\item \href{https://github.com/strands-agents/harness-sdk}{Strands Agents SDK official documentation} (official technology documentation)
\item \href{https://platform.claude.com/docs/en/agent-sdk/overview}{Claude Agent SDK official documentation} (official technology documentation)
\item \href{https://ai-sdk.dev/docs/introduction}{Vercel AI SDK official documentation} (official technology documentation)
\item \href{https://learn.microsoft.com/en-us/agent-framework/overview/}{Microsoft Agent Framework official documentation} (official technology documentation)
\item \href{https://docs.camel-ai.org/}{CAMEL-AI official documentation} (official technology documentation)
\item \href{https://rasa.com/docs/}{Rasa official documentation} (official technology documentation)
\item \href{https://docs.langflow.org/}{Langflow official documentation} (official technology documentation)
\item \href{https://docs.dify.ai/}{Dify official documentation} (official technology documentation)
\item \href{https://docs.flowiseai.com/}{Flowise official documentation} (official technology documentation)
\item \href{https://developers.cloudflare.com/agents/}{Cloudflare Agents SDK official documentation} (official technology documentation)
\item \href{https://huggingface.co/docs/accelerate/index}{Hugging Face Accelerate official documentation} (official technology documentation)
\item \href{https://www.deepspeed.ai/}{DeepSpeed official documentation} (official technology documentation)
\item \href{https://huggingface.co/docs/trl/index}{Hugging Face TRL official documentation} (official technology documentation)
\item \href{https://docs.axolotl.ai/}{Axolotl official documentation} (official technology documentation)
\item \href{https://docs.unsloth.ai/}{Unsloth official documentation} (official technology documentation)
\item \href{https://llm-d.ai/docs/dev}{llm-d official documentation} (official technology documentation)
\item \href{https://llm-d.ai/docs/architecture}{llm-d architecture} (official technology documentation)
\item \href{https://kserve.github.io/website/docs/model-serving/generative-inference/llmisvc/llmisvc-overview}{KServe LLMInferenceService} (official technology documentation)
\item \href{https://docs.nvidia.com/dynamo/latest/}{NVIDIA Dynamo official documentation} (official technology documentation)
\item \href{https://docs.nvidia.com/nim/large-language-models/latest/introduction.html}{NVIDIA NIM official documentation} (official technology documentation)
\item \href{https://docs.nvidia.com/nim/large-language-models/latest/reference/architecture.html}{NIM LLM architecture} (official technology documentation)
\item \href{https://docs.pytorch.org/serve/}{TorchServe official documentation} (official technology documentation)
\item \href{https://deepspeed-mii.readthedocs.io/en/latest/}{DeepSpeed-MII official documentation} (official technology documentation)
\item \href{https://github.com/bentoml/OpenLLM}{OpenLLM official documentation} (official technology documentation)
\item \href{https://lmdeploy.readthedocs.io/en/latest/}{LMDeploy official documentation} (official technology documentation)
\item \href{https://llm.mlc.ai/docs/}{MLC LLM official documentation} (official technology documentation)
\item \href{https://replicate.com/docs}{Replicate official documentation} (official technology documentation)
\item \href{https://docs.runpod.io/serverless/overview}{Runpod Serverless official documentation} (official technology documentation)
\item \href{https://docs.together.ai/docs/inference-overview}{Together AI Inference official documentation} (official technology documentation)
\item \href{https://docs.fireworks.ai/}{Fireworks AI official documentation} (official technology documentation)
\item \href{https://console.groq.com/docs/overview}{GroqCloud official documentation} (official technology documentation)
\item \href{https://inference-docs.cerebras.ai/introduction}{Cerebras Inference official documentation} (official technology documentation)
\item \href{https://docs.aws.amazon.com/bedrock/latest/userguide/what-is-bedrock.html}{Amazon Bedrock official documentation} (official technology documentation)
\item \href{https://developers.cloudflare.com/workers-ai/}{Cloudflare Workers AI official documentation} (official technology documentation)
\item \href{https://docs.databricks.com/aws/en/machine-learning/model-serving/}{Databricks Model Serving official documentation} (official technology documentation)
\item \href{https://www.mongodb.com/docs/atlas/atlas-vector-search/}{MongoDB Atlas Vector Search official documentation} (official technology documentation)
\item \href{https://docs.couchbase.com/server/current/vector-search/vector-search.html}{Couchbase Vector Search official documentation} (official technology documentation)
\item \href{https://docs.databricks.com/aws/en/generative-ai/vector-search/}{Databricks Mosaic AI Vector Search official documentation} (official technology documentation)
\item \href{https://docs.snowflake.com/en/user-guide/snowflake-cortex/cortex-search/cortex-search-overview}{Snowflake Cortex Search official documentation} (official technology documentation)
\item \href{https://docs.oracle.com/en/database/oracle/oracle-database/23/vecse/}{Oracle AI Vector Search official documentation} (official technology documentation)
\item \href{https://docs.tigergraph.com/}{TigerGraph official documentation} (official technology documentation)
\item \href{https://docs.arangodb.com/}{ArangoDB official documentation} (official technology documentation)
\item \href{https://age.apache.org/age-manual/master/intro/overview.html}{Apache AGE official documentation} (official technology documentation)
\item \href{https://docs.janusgraph.org/}{JanusGraph official documentation} (official technology documentation)
\item \href{https://www.comet.com/docs/opik/}{Opik official documentation} (official technology documentation)
\item \href{https://docs.agentops.ai/}{AgentOps official documentation} (official technology documentation)
\item \href{https://docs.parea.ai/}{Parea official documentation} (official technology documentation)
\item \href{https://docs.honeyhive.ai/}{HoneyHive official documentation} (official technology documentation)
\item \href{https://docs.galileo.ai/}{Galileo official documentation} (official technology documentation)
\item \href{https://docs.giskard.ai/}{Giskard official documentation} (official technology documentation)
\item \href{https://inspect.aisi.org.uk/}{Inspect AI official documentation} (official technology documentation)
\item \href{https://www.getmaxim.ai/docs/}{Maxim AI official documentation} (official technology documentation)
\item \href{https://reference.garak.ai/en/latest/}{garak official documentation} (official technology documentation)
\item \href{https://azure.github.io/PyRIT/}{PyRIT official documentation} (official technology documentation)
\item \href{https://github.com/protectai/rebuff}{Rebuff official documentation} (official technology documentation)
\item \href{https://docs.aws.amazon.com/bedrock/latest/userguide/guardrails.html}{Amazon Bedrock Guardrails official documentation} (official technology documentation)
\item \href{https://learn.microsoft.com/en-us/azure/ai-services/content-safety/overview}{Azure AI Content Safety official documentation} (official technology documentation)
\item \href{https://cloud.google.com/security-command-center/docs/model-armor-overview}{Google Cloud Model Armor official documentation} (official technology documentation)
\item \href{https://modelcontextprotocol.io/registry/about}{Official MCP Registry official documentation} (official technology documentation)
\item \href{https://registry.modelcontextprotocol.io/docs}{Registry API} (official technology documentation)
\item \href{https://github.com/github/github-mcp-server}{GitHub MCP Server official documentation} (official technology documentation)
\item \href{https://github.com/awslabs/mcp}{AWS MCP Servers official documentation} (official technology documentation)
\item \href{https://googleapis.github.io/genai-toolbox/}{MCP Toolbox for Databases official documentation} (official technology documentation)
\item \href{https://github.com/microsoft/mcp}{Azure MCP Server official documentation} (official technology documentation)
\item \href{https://docs.zapier.com/mcp/home}{Zapier MCP official documentation} (official technology documentation)
\item \href{https://mcp-use.com/docs/}{mcp-use official documentation} (official technology documentation)
\item \href{https://developers.cloudflare.com/sandbox/}{Cloudflare Sandbox SDK official documentation} (official technology documentation)
\item \href{https://developers.cloudflare.com/sandbox/api/storage/}{Sandbox storage and persistence} (official technology documentation)
\item \href{https://docs.runloop.ai/}{Runloop official documentation} (official technology documentation)
\item \href{https://docs.sprites.dev/}{Sprites official documentation} (official technology documentation)
\item \href{https://docs.browserbase.com/}{Browserbase official documentation} (official technology documentation)
\item \href{https://docs.docker.com/ai/sandboxes/}{Docker Sandboxes official documentation} (official technology documentation)
\item \href{https://docs.wasmtime.dev/}{Wasmtime official documentation} (official technology documentation)
\item \href{https://github.com/google/nsjail}{nsjail official documentation} (official technology documentation)
\item \href{https://argo-workflows.readthedocs.io/}{Argo Workflows official documentation} (official technology documentation)
\item \href{https://docs.flyte.org/}{Flyte official documentation} (official technology documentation)
\item \href{https://docs.metaflow.org/}{Metaflow official documentation} (official technology documentation)
\item \href{https://clear.ml/docs/latest/docs/}{ClearML official documentation} (official technology documentation)
\item \href{https://labelstud.io/guide/}{Label Studio official documentation} (official technology documentation)
\item \href{https://docs.greatexpectations.io/docs/}{Great Expectations official documentation} (official technology documentation)
\item \href{https://docs.airbyte.com/}{Airbyte official documentation} (official technology documentation)
\item \href{https://docs.unstructured.io/}{Unstructured official documentation} (official technology documentation)
\item \href{https://docling-project.github.io/docling/}{Docling official documentation} (official technology documentation)
\end{enumerate}
